\subsection{Fifth letter}

\begin{quotex}
In this letter from \textbf{Rene Guenon} to\textbf{ Guido de Giorgio}, we learn:

\begin{itemize}
\item
Guenon was ignorant of Buddhism, specifically \emph{Vajrayana}\footnote{\url{http://en.wikipedia.org/wiki/Vajrayana}}. 
\item 
Initiation is connected to the Kali Yuga. 
\item
\textbf{Dante}`s symbolism is equivalent to Eastern symbolism. This explains some observations: 

\begin{itemize}
\item
Westerners who don't understand Dante are unlikely to understand Eastern religions. 
\item 
Modern scholars are unlikely to understand Dante. 
\end{itemize}
\item 
Hermetic symbolism is close to Oriental symbolism. 
\item
Historic facts have a symbolic value completely different from what they have in themselves. 
\item
Guenon believes Evola is intelligent but his prejudices often prevent him from seeing things truly. 
\end{itemize}
\end{quotex}

Blois, 15 August, 1927

I am sorry to see from your last letter that your health is still far from being good; I hope however it will improve. Unfortunately, the weather is very bad since we have been here: storms, gales, almost continual rain; it is truly not summer weather. However, here, it typically gets much better, and the sky is not at all the same as in the region of Paris, but this year, the sun is almost always absent. There are neither mountains nor hills, at most some vineyards; but there are forests very close to Blois, and we can go for a number of hikes nearby. This is, I believe, the information that you asked me about; in any case, if you decide to come, don't do it before the first days of September because we are going to be away. I would be very happy to have the chance to finally get acquainted with you, if your health permits you at that time to make the trip without tiring yourself.

\emph{Vajra} is a thunderbolt, but I must warn you that I myself would be currently at pains to say exactly what \emph{Vajrayana} is. I will have to research that on my return to Paris, not having here the necessary information.

I totally agree with your opinion about the current state of humanity in relation to the Kali Yuga. Moreover, it is what I explain in the book that I just completed\footnote{From the date of the letter, this must be \emph{The Crisis of the Modern World}.}. Moreover, it is quite certain that ``initiation" is understood only by the special conditions of the Kali Yuga, outside of which it would not have its \emph{raison d'être}. It is not less true that it is necessary, in fact, to take these conditions into consideration. That is why, while being perfectly in agreement with you in principle, I must however maintain everything that I said about the role of the elite. This role, moreover, is not at all typical of traditions in their religious form. The example of Taoism is a sufficient proof of it, and the same thing is found everywhere, although sometimes in a way less defined (in India, for example). Since the Kali Yuga, ``initiation" exists in the East as well as it used to exist in the West; there it is as a necessity of fact.

Likewise for symbolism: the use of symbols comparable to Hermetic symbols is also totally general, and these symbols are not opposed at all to natural symbols, but, on the contrary, they are connected very normally. Furthermore, the symbolic character of all manifestation permits us to give to historic facts, as well as to all the rest, a value completely different from what they have in themselves. As for what you say about Dante; his symbolism is, if you wish, Western in is exterior form, but it is totally equivalent to Oriental symbolism. Moreover, there arose a true opposition between the East and the West only when the West lost its tradition, including the meaning of symbolism. Hermetism is much closer to the Oriental spirit than to the modern Western spirit. Perhaps we will soon have the opportunity to talk about all of this in more depth.

I received a letter from Evola yesterday, who still insists that I send him something for Ur. I surely believe that I would have to end things by deciding to give him satisfaction. He tells me that he had published the first volume of his Theory of the Absolute Individual, but that he didn't send it to me because of its overly philosophical character, and that he thinks his \textit{Pagan Imperialism}\footnote{\url{https://www.amazon.com/Pagan-Imperialism-Julius-Evola/dp/0999086006/}}, which is due to appear in the autumn, will be liable to interest me more. It is in the same vein of the book I just wrote; I suppose that it is you who spoke to him about it, because I never wrote to anyone else in Italy since I began to prepare it. That is of no importance, since I have already announced it to many other people. I expedited the manuscript to the publisher last week. I am not sorry that this has been completed because until that point, I was unable to take any rest.

It is not Evola, but Reghini, who told me that his attitude toward Catholicism had changed recently. I believe that I explained myself poorly in my last letter, or that perhaps I wrote one name for the other due to a distraction. It is moreover possible, according to what you say, that the same thing is true for both of them. In any case, if you can have any influence on Evola, that would be very felicitous. I believe him to be intelligent, but filled with prejudices of all types. I also think he aspires to a university position and that also can hinder him from many points of view.



\flrightit{Posted on 2012-09-10 by Cologero }

\begin{center}* * *\end{center}

\begin{footnotesize}\begin{sffamily}



\texttt{Avery on 2012-09-11 at 01:09 said: }

``I also think he aspires to a university position and that also can hinder him from many points of view."

I don't know if that guess is correct, but the warning about the limiting nature of the academic mindset is very real. I wonder how Guenon came to this conclusion about Evola.


\hfill

\texttt{Jason-Adam on 2012-09-11 at 01:35 said: }

The remark Guenon made about opposition between East \& West only being a conflict modernity vs Tradition is one to be meditated on. Guenon believes that all traditional civilisations are similar. But do they still conflict with each other ? traditional Catholic Europe contended with traditional Islam. So is it possible there can be a conflict between East \& West even when both are Traditional ? Evola certainly saw the west and east as being different (actions vs contemplation) which Guenon disagreed with.


\hfill

\texttt{Mihai on 2012-09-11 at 09:37 said: }

They are similar in essence, but different in form.

So even though they may have ``internal" unity, the ``external" differences still exist on their own level, sometimes having the aspect of an apparent opposition. It is very normal for different civilizations and traditions to be opposed one another, as different persons are opposed one another, otherwise there would be no history and no manifestation. 

The symbolism of chess may help here. Even tough both armies are, in essence, similar, in order to exist a game they need to be opposed one another.


\end{sffamily}\end{footnotesize}
